% Part: first-order-logic
% Chapter: sequent-calculus
% Section: soundness-identity

\documentclass[../../../include/open-logic-section]{subfiles}

\begin{document}

\olfileid{fol}{seq}{sid}

\olsection{Correctheid met \usetoken{S}{identity}}

\begin{prop}
$\Log{LK}$ met beginsequenten en regels voor identiteit is correct.
\end{prop}

\begin{proof}
Beginsequenten van de vorm ${} \Sequent \eq[t][t]$ zijn geldig, want
voor elke !!{structure}~$\Struct M$ geldt $\Sat{M}{\eq[t][t]}$. (Merk
op dat we veronderstellen dat de term $t$ gesloten is, dat wil zeggen
dat hij geen variabelen bevat, zodat bedelingen niet ter zake doen.)

Stel dat de laatste afleidingsstap in !!a{derivation} een $=$-regel is.
De premisse is dan $\eq[t_1][t_2], \Gamma \Sequent \Delta, !A(t_1)$
en de conclusie is $\eq[t_1][t_2], \Gamma \Sequent \Delta, !A(t_2)$.
Beschouw !!a{structure}~$\Struct M$. We moeten aantonen dat de conclusie
geldig is: als $\Sat{M}{\eq[t_1][t_2]}$ en $\Sat{M}{\Gamma}$, dan geldt
ofwel $\Sat{M}{!C}$ voor een $!C \in \Delta$, ofwel
$\Sat{M}{!A(t_2)}$.

Volgens de inductiehypothese is de premisse geldig. Dit betekent dat,
als $\Sat{M}{\eq[t_1][t_2]}$ en $\Sat{M}{\Gamma}$, ofwel (a) voor een
$!C \in \Delta$ geldt $\Sat{M}{!C}$, ofwel (b)
$\Sat{M}{!A(t_1)}$. In geval (a) zijn we klaar. Beschouw geval (b).
Laat $s$ een bedeling zijn waarvoor $s(x) = \Value{t_1}{M}$. Volgens
\olref[syn][ass]{prop:sentence-sat-true} geldt
$\Sat{M}{!A(t_1)}[s]$. Omdat $\varAssign{s}{s}{x}$, volgt uit
\olref[syn][ext]{prop:ext-formulas} dat $\Sat{M}{!A(x)}[s]$. Uit
$\Sat{M}{\eq[t_1][t_2]}$ volgt $\Value{t_1}{M} = \Value{t_2}{M}$, en
daarom $s(x) = \Value{t_2}{M}$. Door
\olref[syn][ext]{prop:ext-formulas} nogmaals toe te passen, verkrijgen
we ook $\Sat{M}{!A(t_2)}[s]$. Volgens
\olref[syn][ass]{prop:sentence-sat-true} geldt derhalve
$\Sat{M}{!A(t_2)}$.
\end{proof}
\end{document}
